\documentclass[11pt,a4paper]{article}
\usepackage[T1]{fontenc}
\usepackage[utf8]{inputenc}
\usepackage{amsmath,amssymb,amsthm}
\usepackage[margin=1in]{geometry}
\usepackage[colorlinks=true,linkcolor=blue,urlcolor=blue]{hyperref}
\theoremstyle{plain}
\newtheorem{theorem}{Theorem}
\newtheorem{lemma}[theorem]{Lemma}
\newtheorem{proposition}[theorem]{Proposition}
\newtheorem{corollary}[theorem]{Corollary}
\theoremstyle{definition}
\newtheorem{remark}[theorem]{Remark}

\title{A condition on every subblock: the empirical recurrence for OEIS A183848}
\author{Adrian Perez Fontelles\\ \small Independent researcher}
\date{6 September 2026}
\begin{document}
\maketitle

\begin{abstract}
OEIS A183848 carries an empirical recurrence of order $21$ contributed by R.~H.~Hardin. It is
true, and it is decidable rather than empirical. The entry's condition is a condition on every
subblock of a fixed size, so it is decided by a bounded window of consecutive lines; the
admissible windows are the vertices of a finite digraph and the arrays counted are exactly the
walks in it. Hence $a(n)$ is a walk count on $S=26$ vertices and is $C$-finite, and whether the
conjectured recurrence annihilates it is settled by a finite exact computation, with the
Cayley--Hamilton theorem bounding the work.
\end{abstract}

\noindent\small 2020 Mathematics Subject Classification. 05A15, 05B45, 15B36.\normalsize

\section{The conjecture}

OEIS A183848 is ``Number of (n+1)X4 0..3 arrays with each element of every 2X2 subblock being the sum mod 4 of two others.''. It has offset $1$ and begins
\[
328, 1766, 9557, 54724, 310856, 1846987, 10793131, 65461226,\ \dots
\]
The entry states, as a conjecture contributed by its author and never marked settled:
\begin{quote}
Empirical: a(n)=8*a(n-1)+59*a(n-2)-565*a(n-3)-822*a(n-4)+13359*a(n-5)-4545*a(n-6)-132397*a(n-7)+154942*a(n-8)+558160*a(n-9)-920421*a(n-10)-1048403*a(n-11)+2303548*a(n-12)+729071*a(n-13)-2730029*a(n-14)+118317*a(n-15)+1498861*a(n-16)-308470*a(n-17)-342440*a(n-18)+95112*a(n-19)+25872*a(n-20)-8064*a(n-21)
\end{quote}
As of the ``Last modified'' line on the live entry (Oct 03 2025 02:06:48, revision 9) this is still
recorded as empirical, and nothing on the entry records it as proved.

\section{The count is a walk count}

The entry counts arrays with $4$ columns over the alphabet $\{0,1,2,3\}$, growing one row at a time, and the condition it imposes is that in every $2\times2$ subblock, each of the four entries must be congruent modulo $4$ to the sum of two of the other three.

That condition is decided by a bounded window. A $2\times2$ subblock occupies $2$ consecutive rows, so a window of $2$ consecutive rows decides everything that the arrival of a new row settles; nothing outside the window is consulted. Take those windows as the vertices of a digraph, with an edge from $u$ to $v$ when $v$'s new row may follow $u$, and merge vertices with identical futures; that leaves $S=26$. An admissible array is exactly a walk, so with $M$ the adjacency matrix and $\iota,\tau$ the vectors recording which windows may begin and end an array,
\[
a(n)\;=\;\iota^{\!\top}M^{\,n}\tau ,
\]
the exponent fixed by the entry's own shape and checked against its published terms. In
particular $a$ is $C$-finite.

\section{The proof}

\begin{lemma}\label{lem:crit}
Let $q(t)=t^{21}-\sum_i c_it^{21-i}$ be the characteristic polynomial of the
conjectured recurrence. The residual $a(n)-\sum_ic_ia(n-i)$ equals
$\iota^{\!\top}M^{\,j}q(M)\tau$ for the corresponding walk index $j$, so the recurrence
holds from some point on if and only if $\iota^{\!\top}M^{j}q(M)\tau=0$ for all large $j$,
and it is enough to examine $j<S$.
\end{lemma}

\begin{proof}
Factor $M^{j}$ out of $M^{j+21}-\sum_ic_iM^{j+21-i}$. For the bound, $M^{S}$
is an integer combination of $I,M,\dots,M^{S-1}$ by Cayley--Hamilton, so the numbers
$u_j=\iota^{\!\top}M^{j}q(M)\tau$ obey the monic recurrence given by the characteristic
polynomial of $M$; $S$ consecutive zeros force every later one, and the last nonzero $u_j$
pins the threshold exactly. A constant factor in front of $a$ divides out of the residual and
changes nothing here.
\end{proof}

\begin{theorem}
\[
a(n)\;=\;8\,a(n-1) + 59\,a(n-2) - 565\,a(n-3) - 822\,a(n-4) + 13359\,a(n-5) - 4545\,a(n-6) - 132397\,a(n-7) + 154942\,a(n-8) + 558160\,a(n-9) - 920421\,a(n-10) - 1048403\,a(n-11) + 2303548\,a(n-12) + 729071\,a(n-13) - 2730029\,a(n-14) + 118317\,a(n-15) + 1498861\,a(n-16) - 308470\,a(n-17) - 342440\,a(n-18) + 95112\,a(n-19) + 25872\,a(n-20) - 8064\,a(n-21)
\]
for every $n>21$.
\end{theorem}

\begin{proof}
The vector $w=q(M)\tau$ was formed by $21$ matrix--vector products in exact integer
arithmetic and $\iota^{\!\top}M^{j}w$ evaluated for $j=0,1,\dots$, again exactly; those
integers vanish from the index corresponding to $n=21+1$ onwards and the run continued
until the working vector was identically zero, which settles every larger $j$ at once.
\end{proof}

\section{Verification}

Four checks, each independent of the algebra above.

First, the model reproduces all $18$ terms the entry publishes, exactly, in integer
arithmetic. This is what ties the model to the entry: the digraph is built by reading the
entry's English, and a misreading gives different counts.

Second, the reading was pinned before the digraph was written, by a program that enumerates
every array of the given shape one by one and tests the entry's condition on each subblock
directly. That program shares no code with the transfer matrix, so an error in one does not
hide an error in the other.

Third, the conjectured recurrence was evaluated directly on the published terms, with no
matrices involved, and holds wherever the proved range and the data overlap.

Fourth, the annihilation test of Lemma~\ref{lem:crit} was carried out in exact integer
arithmetic on the whole vector rather than on a sampled prefix.

\begin{thebibliography}{9}
\bibitem{oeis} The OEIS Foundation, \emph{The On-Line Encyclopedia of Integer
Sequences}, \url{https://oeis.org}, 2026. Sequence A183848.
\bibitem{stanley} R.~P.~Stanley, \emph{Enumerative Combinatorics}, Volume 1, 2nd ed.,
Cambridge University Press, 2011. (Transfer-matrix method, Section 4.7.)
\bibitem{fs} P.~Flajolet and R.~Sedgewick, \emph{Analytic Combinatorics}, Cambridge
University Press, 2009.
\bibitem{horn} R.~A.~Horn and C.~R.~Johnson, \emph{Matrix Analysis}, 2nd ed., Cambridge
University Press, 2013.
\end{thebibliography}

\end{document}
